\chapter{Deep Dive: The Browser Meets the Gateway --- Token Storage, Refresh Flow, and CORS}
\label{ch:dd-frontend-auth}

Chapter~26 put the React application on the same origin as the API and
noted, almost in passing, that CORS ``is simply never consulted''.
Chapter~28 pushed the token design from the server's side. This deep
dive takes the other end of the wire: the browser, which is the one
runtime in the system that you do not control, that runs code injected
by anyone who can get a script tag onto your page, and that will
happily send a cookie to an attacker's form. Every decision about where
a token lives is a decision about which of those threats you accept.

Three questions organize the chapter:

\begin{enumerate}
  \item \textbf{Where does the token live} in the browser, and which
  attack does each home expose it to?
  \item \textbf{What happens at expiry} --- how the refresh flow runs,
  and why ten tabs must produce one refresh call?
  \item \textbf{What is CORS protecting}, what does it cost per
  request, and why does one origin make it vanish?
\end{enumerate}

\section{Why this, and what it solves}

The server issues two credentials (Chapter~5): a 15-minute JWT access
token and a 7-day opaque refresh token. The browser must keep both
across page reloads, attach the first to every API call, and exchange
the second for a fresh first when it expires --- without the user
noticing. The project's answer is a Zustand store with the
\code{persist} middleware and an axios interceptor. Read the store
first:

\begin{lstlisting}
export const useAuthStore = create<AuthState>()(
  persist(                                            (*@\co{1}@*)
    (set) => ({
      accessToken: null,
      refreshToken: null,
      user: null,
      isAuthenticated: false,
      setTokens: (accessToken, refreshToken) =>
        set({ accessToken, refreshToken, isAuthenticated: true }),
      setUser: (user) => set({ user }),
      logout: () =>
        set({ accessToken: null, refreshToken: null,
              user: null, isAuthenticated: false }),
    }),
    { name: 'auth-storage' }                          (*@\co{2}@*)
  )
);
\end{lstlisting}

\begin{description}
  \item[\co{1}] \code{persist} writes the whole state to storage on
  every change and rehydrates it on page load. That is what survives a
  reload; without it every refresh of the page would be a logout.
  \item[\co{2}] The storage key. The default backend is
  \code{localStorage}, so open DevTools, Application, Local Storage,
  and there they are: both tokens, in clear text, under
  \code{auth-storage}. Any script running on the page can read them
  with one line.
\end{description}

\begin{decision}[localStorage, for a first version]
The project stores both tokens in \code{localStorage} because it is the
simplest thing that survives a reload, works identically on
\code{localhost:3000} and \code{ts.local}, needs no cookie
configuration on the server, and sidesteps CSRF entirely --- a token
that JavaScript attaches by hand is never sent by the browser on its
own. The price is that the tokens are exactly as safe as every script
on the page. That was acceptable while the page loads only its own
bundle. It stops being acceptable the day a third-party script (an
analytics tag, a chat widget, a compromised npm dependency) runs on the
same origin: any of them can read \code{auth-storage} and ship a 7-day
refresh token elsewhere. The upgrade in Section~\ref{sec:fe-enhance}
moves the refresh token into an \code{HttpOnly} cookie and keeps the
access token in memory --- more moving parts, but a script can no
longer steal the long-lived credential.
\end{decision}

\section{The refresh flow, as written}

\subsection{The interceptor}

Every API call goes through one axios instance:

\begin{lstlisting}
const apiClient = axios.create({ baseURL: '/api' });   (*@\co{1}@*)

apiClient.interceptors.request.use((config) => {
  const token = useAuthStore.getState().accessToken;
  if (token) config.headers.Authorization = `Bearer ${token}`;
  return config;                                        (*@\co{2}@*)
});

apiClient.interceptors.response.use(
  (response) => response,
  async (error) => {
    const originalRequest = error.config;
    if (error.response?.status === 401
        && !originalRequest._retry) {                   (*@\co{3}@*)
      originalRequest._retry = true;
      const refreshToken = useAuthStore.getState().refreshToken;
      if (refreshToken) {
        try {
          const { data } = await axios.post(            (*@\co{4}@*)
              '/api/auth/refresh', { refreshToken });
          useAuthStore.getState()
              .setTokens(data.accessToken, data.refreshToken);
          originalRequest.headers.Authorization =
              `Bearer ${data.accessToken}`;
          return apiClient(originalRequest);            (*@\co{5}@*)
        } catch {
          useAuthStore.getState().logout();
          window.location.href = '/login';
        }
      }
    }
    return Promise.reject(error);
  }
);
\end{lstlisting}

\begin{description}
  \item[\co{1}] A relative base URL. In development Vite's proxy
  forwards \code{/api} to \code{localhost:8080}; in the cluster Traefik
  does (Chapter~26). The browser never sees a second origin.
  \item[\co{2}] The bearer header is attached by code on every request.
  This is the CSRF immunity mentioned above: a form on
  \code{evil.example} cannot make the browser add this header.
  \item[\co{3}] A 401 from the gateway --- the JSON body ``Invalid or
  expired token'' from \code{JwtAuthenticationFilter} --- triggers one
  retry per request, flagged on the config object so a second 401
  falls through to the caller.
  \item[\co{4}] The refresh call deliberately uses the \emph{bare}
  \code{axios}, not \code{apiClient}: otherwise a 401 from
  \code{/refresh} itself would re-enter this interceptor and loop.
  \item[\co{5}] The original request is replayed with the new token.
  The caller's \code{await} resolves as if nothing happened.
\end{description}

This is correct for one request. It is wrong for a page.

\subsection{The single-flight problem}

\code{DashboardPage} mounts and TanStack Query fires its queries in
parallel: courses, assignments, the user profile, dictionary counts.
If the access token expired while the laptop was closed, all of them
return 401 within the same few milliseconds, and each one runs the
interceptor independently. Four requests, four calls to
\code{/api/auth/refresh}, four new access tokens, and --- because
\code{setTokens} runs four times --- whichever response lands last
wins. Today that is harmless because auth-service returns the same
refresh token every time (Chapter~28 called this out as the missing
rotation). Add rotation, which you should, and the race becomes fatal:
the first refresh invalidates the token the other three are about to
present, they fail, and the user is logged out by the very mechanism
meant to keep them in.

The fix is a \emph{single-flight} refresh: the first 401 starts one
refresh and stores its promise; every concurrent 401 awaits the same
promise instead of starting its own; when it resolves, all of them
replay with the same new token. Section~\ref{sec:fe-enhance} has the
complete code. Interviewers ask this as ``ten tabs hit 401 at once ---
how many refresh calls?'' and the senior answer is ``one per tab
without coordination, one in total with a lock across tabs'', followed
by naming the tool for the cross-tab case: the \code{Web Locks} API or
a \code{BroadcastChannel}, since each tab has its own JavaScript heap
and cannot see another tab's pending promise.

\subsection{Logout and the OAuth callback}

\code{Header.handleLogout} posts the refresh token to
\code{/api/auth/logout}, which deletes its row, then clears the store
and navigates to \code{/login}. That revokes the 7-day credential
immediately. The 15-minute access token is not revoked --- it cannot
be, per Chapter~28 --- so a copy taken before logout keeps working
until it expires. ``How do you log a user out everywhere?'' has the
same answer with one more step: delete \emph{all} of the user's
refresh-token rows, not just the one this browser holds, and accept
the 15-minute window or add the \code{jti} deny-list.

The Google login lands on \code{OAuthCallbackPage} with both tokens as
query parameters --- \code{OAuth2LoginSuccessHandler} builds the
redirect that way. The page reads them, stores them, and navigates
with \code{replace: true} so the URL with the tokens does not stay in
history. That last detail matters more than it looks: a URL is copied
into the browser history, into the \code{Referer} header of any
resource the page loads next, and into server access logs. The
\code{replace} shortens the exposure; it does not remove it from the
gateway's and auth-service's access logs, which now contain a valid
refresh token in plain text.

\section{CORS: the policy you are not using}

Same-origin policy is the browser's rule that a script loaded from
origin A may not \emph{read} responses from origin B. CORS is the
protocol by which B can relax that rule per origin. The project's
gateway has a \code{CorsConfig} that lists \code{localhost:3000} and
\code{localhost:5173}, allows all methods and headers, and sets
\code{allowCredentials(true)}. Read it against the two ways the app is
actually served:

\begin{itemize}
  \item \textbf{Development.} Vite serves the page on \code{:3000} and
  \emph{proxies} \code{/api} to \code{:8080}. The browser sends
  requests to \code{:3000}; Vite forwards them server-side. No
  cross-origin request exists, so no preflight, so \code{CorsConfig}
  is never exercised.
  \item \textbf{Cluster.} Traefik routes \code{/api} to the gateway
  and everything else to nginx on the same host (Chapter~26). Same
  origin again.
\end{itemize}

\code{CorsConfig} therefore serves exactly one scenario: someone
running the built bundle from a different origin --- say, opening
\code{dist/index.html} through a static server on \code{:5173} without
the proxy. It is a safety net, and a permissive one:
\code{addAllowedHeader("*")} plus credentials is the widest policy
Spring will accept without throwing.

What CORS would cost if it were active is worth knowing for the
interview. A request with a JSON body and an \code{Authorization}
header is not a ``simple request'', so the browser first sends an
\code{OPTIONS} preflight asking whether the method and headers are
allowed. That is one extra round trip per distinct request, cacheable
by \code{Access-Control-Max-Age}, which the project's config does not
set --- so every call would preflight. On a 40\,ms link, every API
call becomes 80\,ms before any work happens. One origin removes that
cost entirely, which is the performance argument for Chapter~26's
design on top of the security one.

\section{Pros and cons}

\begin{center}
\begin{tabular}{@{}p{0.30\linewidth}p{0.32\linewidth}p{0.32\linewidth}@{}}
\toprule
 & \textbf{localStorage + bearer (today)} & \textbf{HttpOnly cookie} \\
\midrule
XSS steals the token & yes, both tokens, one line of JS & no --- script
cannot read the cookie \\
CSRF & immune: header is attached by code & exposed unless
\code{SameSite} + CSRF token or custom-header check \\
Survives reload & yes, via \code{persist} & yes, the browser keeps it \\
Multi-tab & each tab reads storage; refresh races & same cookie in every
tab; refresh races still \\
Works cross-origin & yes & needs \code{Allow-Credentials} and an exact
origin list \\
Server changes & none & set-cookie on login, read-cookie on refresh \\
\bottomrule
\end{tabular}
\end{center}

The honest summary a senior engineer gives: bearer-in-storage trades a
CSRF problem you can solve for an XSS problem you cannot. Any injected
script defeats it completely. Cookies trade the other way, and their
problem has known, mechanical fixes.

\section{What breaks}

\begin{pitfall}[four refreshes, one winner]
Symptom: after leaving the app open overnight, the first click works,
then the user is sent to \code{/login} even though nothing is wrong
with their account. Cause: parallel 401s each ran the interceptor;
with refresh-token rotation on the server, the second refresh
presented a token the first had just invalidated, and the
\code{catch} branch logged the user out. Fix: the single-flight
interceptor in Section~\ref{sec:fe-enhance}. Without rotation the
symptom is only wasted calls --- which is why the race is invisible
today.
\end{pitfall}

\begin{pitfall}[the token in the URL]
Symptom: a refresh token appears in the gateway's access log, in
Traefik's log, and in the browser history of a shared computer. Cause:
\code{OAuth2LoginSuccessHandler} passes both tokens as query
parameters of the redirect to \code{/oauth/callback}. Fix: put the
tokens in the URL \emph{fragment} (\code{\#accessToken=...}) --- the
fragment is never sent to any server --- or better, return only a
one-time code that the page exchanges for tokens with a POST. The
cookie design below removes the refresh token from the URL entirely.
\end{pitfall}

\begin{pitfall}[a 401 that is not about the token]
Symptom: a wrong password on the login page briefly flashes the app
shell, then returns to \code{/login}. Cause: auth-service answers a
bad login with 401; the interceptor treats \emph{every} 401 as ``token
expired'' and, if a stale refresh token exists in storage from an
earlier session, tries to refresh, fails, and runs \code{logout()} plus
a hard redirect. Fix: exempt the auth endpoints from the refresh path
(\code{originalRequest.url} starting with \code{/auth/}), and clear the
store on the login page's mount.
\end{pitfall}

\section{What breaks at 3\,a.m.}

Three failures that only the client sees.

\begin{itemize}
  \item \textbf{The laptop closed for eight days.} On return the
  access token is long dead and the refresh token expired yesterday.
  The first request gets 401, the refresh gets 401 (auth-service
  deletes the expired row and throws), the interceptor logs the user
  out. Correct behaviour --- but every in-flight query on the page
  fails first, and TanStack Query may retry them three times each
  before the redirect lands. Set \code{retry: false} for 401s in the
  query client's default options, or the user watches spinners for
  ten seconds before being sent to login.
  \item \textbf{Client clock skew.} The project never checks
  \code{exp} on the client, and that is right: the client's clock is
  not trustworthy and the gateway decides. A ``smart'' client that
  pre-emptively refreshes when \code{exp} is near will, on a machine
  whose clock is 20 minutes fast, refresh on every request --- or, 20
  minutes slow, never refresh and get 401 anyway. Let the server's 401
  be the signal.
  \item \textbf{CORS misconfigured in production only.} The symptom is
  a console error, ``No Access-Control-Allow-Origin header'', on a
  request that returns 200 in the network tab. The browser got the
  response and refused to hand it to the script. It happens when the
  page and API are on different origins in production (a CDN for the
  bundle, the gateway on another host) and the gateway's allowed list
  still says \code{localhost}. The check is always the same: compare
  the \code{Origin} request header with the \code{allowedOrigins}
  list, character for character, scheme and port included. And note
  the shape of the error: it is never visible to the server, whose
  logs show a successful request.
\end{itemize}

\section{Enhancing the resolution}
\label{sec:fe-enhance}

The target design: the access token lives only in JavaScript memory
(lost on reload, restored by a silent refresh); the refresh token lives
in an \code{HttpOnly} cookie that JavaScript cannot read and the
browser sends only to the refresh endpoint.

\textbf{1. Auth-service sets the cookie.} On login, 2FA verification,
and refresh, the controller returns the access token in the body and
the refresh token in a cookie:

\begin{lstlisting}[language=Java]
private ResponseCookie refreshCookie(String value, long maxAgeSec) {
    return ResponseCookie.from("refresh_token", value)
            .httpOnly(true)       // no document.cookie access
            .secure(true)         // HTTPS only; false on plain localhost
            .sameSite("Strict")   // never sent from another site
            .path("/api/auth/refresh")   // only this endpoint sees it
            .maxAge(maxAgeSec)
            .build();
}

@PostMapping("/login")
public ResponseEntity<LoginResponse> login(
        @Valid @RequestBody LoginRequest request) {
    LoginResponse r = authService.login(request);
    if (r.refreshToken() == null) {          // 2FA or password change
        return ResponseEntity.ok(r);
    }
    return ResponseEntity.ok()
            .header(HttpHeaders.SET_COOKIE, refreshCookie(
                    r.refreshToken(),
                    jwtService.getRefreshTokenExpiryMs() / 1000).toString())
            .body(new LoginResponse(r.accessToken(), null,
                    false, false, null));    // token not in the body
}

@PostMapping("/refresh")
public ResponseEntity<LoginResponse> refresh(
        @CookieValue(name = "refresh_token", required = false)
        String cookie) {
    if (cookie == null) {
        return ResponseEntity.status(HttpStatus.UNAUTHORIZED).build();
    }
    LoginResponse r = authService.refreshToken(
            new TokenRefreshRequest(cookie));
    return ResponseEntity.ok()
            .header(HttpHeaders.SET_COOKIE, refreshCookie(
                    r.refreshToken(),
                    jwtService.getRefreshTokenExpiryMs() / 1000).toString())
            .body(new LoginResponse(r.accessToken(), null,
                    false, false, null));
}

@PostMapping("/logout")
public ResponseEntity<Void> logout(
        @CookieValue(name = "refresh_token", required = false)
        String cookie) {
    if (cookie != null) authService.logout(cookie);
    return ResponseEntity.noContent()
            .header(HttpHeaders.SET_COOKIE,
                    refreshCookie("", 0).toString())   // delete it
            .build();
}
\end{lstlisting}

The cookie's \code{path} is the strongest of the four attributes: even
if some other endpoint were vulnerable, the browser never attaches this
cookie to it. Note the path is the \emph{gateway's} view
(\code{/api/auth/refresh}), because the browser talks to the gateway;
the gateway strips the prefix afterwards. \code{SameSite=Strict} is
the CSRF defence: a cross-site form POST does not carry the cookie at
all. The JSON body requirement on \code{/refresh} is a second layer ---
a plain HTML form cannot send \code{application/json} without a
preflight, which the same-origin design would not grant to another
site.

\textbf{2. The client keeps the access token in memory.} The store
drops \code{persist} and the refresh token entirely:

\begin{lstlisting}
export const useAuthStore = create<AuthState>()((set) => ({
  accessToken: null,          // memory only: gone on reload, by design
  user: null,
  isAuthenticated: false,
  setAccessToken: (accessToken) =>
    set({ accessToken, isAuthenticated: true }),
  setUser: (user) => set({ user }),
  logout: () =>
    set({ accessToken: null, user: null, isAuthenticated: false }),
}));
\end{lstlisting}

On page load nothing is in memory; the app performs one silent refresh
before rendering protected routes. A reload therefore costs one
round-trip to \code{/api/auth/refresh}, and if the cookie is gone the
user is at \code{/login} --- the same behaviour as an expired session.

\textbf{3. Single-flight refresh with replay.} The interceptor keeps
one pending promise and a queue of requests waiting on it:

\begin{lstlisting}
let refreshing: Promise<string> | null = null;   // module-level lock

async function refreshAccessToken(): Promise<string> {
  if (!refreshing) {                              // first 401 starts it
    refreshing = axios
      .post('/api/auth/refresh', null, { withCredentials: true })
      .then(({ data }) => {
        useAuthStore.getState().setAccessToken(data.accessToken);
        return data.accessToken as string;
      })
      .finally(() => { refreshing = null; });     // release the lock
  }
  return refreshing;                              // others await it
}

apiClient.interceptors.response.use(
  (response) => response,
  async (error) => {
    const original = error.config;
    const isAuthCall = original.url?.startsWith('/auth/');
    if (error.response?.status !== 401 || original._retry
        || isAuthCall) {
      return Promise.reject(error);
    }
    original._retry = true;
    try {
      const token = await refreshAccessToken();    // shared promise
      original.headers.Authorization = `Bearer ${token}`;
      return apiClient(original);                  // replay
    } catch {
      useAuthStore.getState().logout();
      window.location.href = '/login';
      return Promise.reject(error);
    }
  }
);
\end{lstlisting}

Ten concurrent 401s now produce one POST; the other nine
\code{await} the same promise and replay with the token it resolves
to. \code{withCredentials: true} is what makes axios send the cookie
--- and what makes the request \emph{credentialed} for CORS purposes.
For cross-tab coordination, wrap the body of
\code{refreshAccessToken} in \code{navigator.locks.request('refresh',
...)}: the Web Locks API is process-wide, so tabs take turns, and the
second tab's refresh finds a rotated cookie already set by the first
and simply succeeds.

\textbf{4. The gateway, if a second origin ever appears.} Credentialed
requests forbid the wildcard; the origin must be listed exactly and the
preflight cached:

\begin{lstlisting}[language=yaml]
spring:
  cloud:
    gateway:
      globalcors:
        cors-configurations:
          '[/**]':
            allowed-origins:
              - https://app.teachersupporter.example   # exact, with scheme
            allowed-methods: [GET, POST, PUT, PATCH, DELETE, OPTIONS]
            allowed-headers: [Authorization, Content-Type]
            allow-credentials: true                   # cookies cross-origin
            max-age: 3600                             # cache the preflight
\end{lstlisting}

With \code{SameSite=Strict} the cookie would still not travel to a
different site, so a genuinely cross-site deployment needs
\code{SameSite=None; Secure} and accepts the CSRF exposure that comes
with it --- the reason the same-origin design of Chapter~26 is the one
to keep.

\begin{handson}[watch the cookie do its job]
After deploying the cookie version, log in through the UI, then inspect
what the browser holds and what it sends:
\begin{lstlisting}[language=bash]
$ curl -si -c jar.txt -X POST http://localhost:8080/api/auth/login \
    -H 'Content-Type: application/json' \
    -d '{"email":"teacher@ts.local","password":"secret"}' | grep -i set-cookie
Set-Cookie: refresh_token=3f9c...; Path=/api/auth/refresh; Max-Age=604800;
  HttpOnly; SameSite=Strict
$ curl -s -b jar.txt -X POST http://localhost:8080/api/auth/refresh \
    | jq .accessToken
"eyJhbGciOiJIUzI1NiJ9..."
$ curl -s -b jar.txt -o /dev/null -w '%{http_code}\n' \
    http://localhost:8080/api/courses
401
\end{lstlisting}
The last line is the \code{path} attribute at work: the cookie exists
in the jar but curl, like a browser, does not send it to
\code{/api/courses}. In DevTools the same cookie shows a padlock in the
HttpOnly column and \code{document.cookie} returns an empty string.
\end{handson}

\section{Outside the project}

The pattern above --- access token in memory, refresh token in an
\code{HttpOnly} cookie, silent refresh on load --- is what Auth0's
SPA SDK, Okta's, and most bank web apps implement, usually with refresh
rotation and reuse detection on the server. The alternative that
removes the browser from token handling altogether is the
\emph{backend-for-frontend}: a small server (or the gateway itself)
holds the tokens in a server-side session, and the browser carries only
an ordinary session cookie; Spring Cloud Gateway supports this with its
\code{TokenRelay} filter and an OIDC login. Interviewers ask ``where
do you store the JWT in the browser?'' hoping for ``nowhere that a
script can read it, and the refresh token never in the URL'' followed
by the trade-off table above. ``What about XSS?'' has one answer that
is not about storage at all: a Content Security Policy that stops
foreign scripts from running in the first place, which is the defence
every storage choice still depends on.

\section{Questions from real interviews}

\subsection*{Q: Where do you store the JWT in the browser, and why?}

Start by separating the two tokens, because they have different
lifetimes and different blast radii. The 15-minute access token can
live in JavaScript memory: losing it on reload costs one silent refresh,
and a script that steals it has fifteen minutes. The refresh token is
the credential that matters --- seven days here --- and it belongs in an
\code{HttpOnly}, \code{Secure}, \code{SameSite=Strict} cookie scoped by
\code{Path} to the refresh endpoint, so no script can read it and the
browser only ever sends it to one URL. \code{localStorage}, which this
project uses today, is defensible for a first version and indefensible
once any third-party script runs on the page: one \code{JSON.parse}
away from a week of access. The senior close is that storage is the
second line; the first is a Content Security Policy that keeps
injected scripts from executing at all.

\subsection*{Q: Ten tabs get a 401 at the same moment. How many refresh calls happen?}

With the interceptor as written, ten --- each tab has its own memory
and its own pending request, and within a single tab each parallel
query runs the interceptor separately, so it can be more than ten. With
refresh-token rotation on the server that is not waste, it is an
outage: the first refresh rotates the token, the other nine present the
old one and are rejected, and the reuse detection that good servers
implement revokes the whole family. Inside one tab the fix is a shared
promise --- first 401 starts the refresh, later 401s await it, all
replay with the result. Across tabs it is \code{navigator.locks}: tabs
take the lock in turn, and the second finds the cookie already rotated
and succeeds with a normal refresh. The follow-up worth volunteering is
that the server should tolerate a short grace window on the previous
token, a few seconds, because network reordering means the ``old''
token can legitimately arrive after the rotation.

\subsection*{Q: XSS versus CSRF --- which one does each storage choice expose you to?}

Bearer tokens in storage: immune to CSRF, because the browser never
attaches an \code{Authorization} header on its own, and fully exposed
to XSS, because any script on the origin reads storage. Cookies: immune
to XSS reading when \code{HttpOnly}, and exposed to CSRF because the
browser attaches cookies to any request to that host, including one
from an attacker's page --- unless \code{SameSite=Strict} or
\code{Lax} stops it, and a custom-header or token check backs that up
for older clients. So the real choice is between a problem with a
mechanical fix (CSRF) and one without (XSS). A senior engineer adds
that XSS on a page with an \code{HttpOnly} cookie can still \emph{use}
the session while the page is open --- it just cannot exfiltrate the
credential --- which is why short access tokens and server-side
revocation still matter.

\subsection*{Q: A user closes the laptop for eight days and comes back. What happens?}

Both credentials are dead: the access token by fifteen minutes, the
refresh token by a day. The first API call gets 401, the interceptor
calls refresh, auth-service finds the row expired, deletes it and
answers 401, and the client logs out and redirects to
\code{/login}. That is the correct outcome; the quality question is
what the user sees on the way. Every query that was in flight fails
first, and if the query client retries failed requests three times
with backoff --- TanStack Query's default --- the redirect is preceded
by several seconds of spinners and repeated 401s in the log. Configure
the query client not to retry on 401, run one refresh on app mount
before rendering protected routes, and the eight-day return is a clean
bounce to login with the intended route preserved for after.

\subsection*{Q: How do you log a user out everywhere?}

The refresh side is a database operation: delete every row in
\code{refresh\_tokens} for that user, not only the one the current
browser presented. Every other device fails its next refresh and lands
on login within fifteen minutes. The access side is the Chapter~28
problem --- a signed token cannot be recalled --- so state the window
honestly and name the two ways to close it: a \code{jti} deny-list in
Redis checked by the gateway, or shorter access tokens. In this project
the concrete change is one repository method,
\code{deleteAllByUser(user)}, called from a new
\code{/auth/logout-all} endpoint and from the admin's disable-user
action. Mention the interview trap: a client-side ``logout'' that only
clears storage revokes nothing --- the token is still valid in whoever
copied it.

\subsection*{Q: A CORS error appears in production but never in development. What do you check?}

First, whether the request is cross-origin at all in production and not
in development --- in this project Vite's proxy and Traefik both keep
the page and the API on one origin, so a CORS error in production means
the deployment split them: a CDN for the bundle, a different host or
port for the gateway. Then compare the \code{Origin} header the browser
sent, byte for byte including scheme and port, against the gateway's
allowed list; \code{https://app.example} and
\code{https://app.example:443} are different strings to Spring. Then
the credentials rule: with cookies or \code{withCredentials}, the
wildcard origin is rejected by the browser even if the server sends it,
and \code{Access-Control-Allow-Credentials: true} must be present. Then
the preflight: an \code{OPTIONS} that hits a route which strips the
prefix and forwards to a service returns that service's 401 instead of
the CORS headers, so CORS must be answered at the gateway before
routing --- which \code{CorsWebFilter} does. Last, look at the server
logs and expect to find nothing: the request succeeded there; the
browser is the one refusing, and that is the shape of every CORS
problem.

\begin{quiz}
\begin{enumerate}
  \item Open DevTools on the running app and name the storage key that
  holds both tokens. Write the one-line JavaScript expression an
  injected script would use to read the refresh token, and say what it
  can do with it for the next seven days.
  \item The interceptor calls bare \code{axios} for the refresh
  instead of \code{apiClient}. Explain the loop that would otherwise
  occur, step by step, starting from a 401 on \code{/refresh}.
  \item With refresh-token rotation enabled on the server, trace four
  parallel 401s through the interceptor as written and state at which
  step the user is logged out. Then explain why the single-flight
  version avoids it.
  \item The cookie design sets \code{Path=/api/auth/refresh}. Give two
  distinct attacks this attribute blunts, and explain why the path is
  written from the gateway's point of view rather than
  auth-service's.
  \item \code{CorsConfig} lists two origins and allows credentials. In
  which of the project's two deployment modes is it consulted, and what
  single change to the deployment would make it matter?
  \item A user reports a CORS error; the gateway log shows the request
  returned 200. Explain why both observations are true at once, and
  list the three header values you would compare first.
\end{enumerate}
\end{quiz}
